% Section S12, second subsection: full proofs of the two statements of Section S2 whose proofs are not in S12.1
% (the theorem on outcomes with spare moves, and the local test for (V2)). Items are numbered A.18, A.19 by hand.

This subsection proves two statements of Section~\ref{S-kings}: Theorem~\ref{thm:loopy} on outcomes with spare moves
(Theorem~A.18a), and the frozen-king test for (V2) (Proposition~A.19 and Corollary~A.19a). They are Theorem~3.4 and
Proposition~2.11 with Corollary~2.11a of the working note \texttt{kings/theory/LOOPY.md} in the code copy. We use the
notation of Parts A.I--A.III and cite items A.1--A.17 of the previous subsection. Every further definition is given
here, and every proof is a hand proof. The one computed statement (in Remark~A.19b) is marked as such and is used
nowhere.

\subsubsection*{A.V Games with decisive terminals and spare moves}

\textbf{D-games.} A \emph{D-game} $G$ is a finite set of positions with, for each position $x$, a finite set $L(x)$ of
White options and a finite set $R(x)$ of Black options. An option is a position or one of two \emph{terminals}:
$\oplus$, ``White has won'', and $\ominus$, ``Black has won''. A move to a terminal is a \emph{decisive event}; a move of
$\sigma$ to the terminal of $\bar\sigma$ is a \emph{suicide}. Play from a position $x$ with a side $s$ to move
alternates, the mover choosing one of its options. A move to a terminal ends play with that winner, a player with no
option loses, and infinite play is a draw. A side $\sigma$ \emph{wins} $(x,s)$ if it has a strategy that ends every play
with a win for $\sigma$ after finitely many moves. At most one side wins (play the two strategies against each other);
if neither does, $(x,s)$ is a \emph{draw}. $G$ is \emph{short} if the directed graph on its positions, with an edge from
$x$ to every position in $L(x)\cup R(x)$, is acyclic; a move from $x$ to $x$ is a cycle. Then a sequence of moves
between positions, made by either side in any order, visits distinct positions, so it has at most $h-1$ moves, where $h$
is the number of positions of $G$; and the positions can be listed so that every option of a position that is itself a
position comes before it.

\textbf{Sums.} The sum $A_1\boxplus\dots\boxplus A_k$ of D-games has the tuples of positions as positions. A move is a
move in one coordinate, and a move to a terminal in any coordinate ends the whole play with that terminal: every
decisive event ends the game. The sum is a D-game (an option of a tuple is a terminal, or the tuple with one coordinate
replaced by an option of it), and it is short if every $A_i$ is, since a move between tuples moves one coordinate
to an earlier position in its listing and leaves the others fixed.

\textbf{Spare moves.} The king group $Z$ of Section~\ref{S-kings}, in which a side that can wait passes and every other
move loses, is modelled by a D-game, also called $Z$, with a single position $z$ and $L(z)\subseteq\{z,\ominus\}$,
$R(z)\subseteq\{z,\oplus\}$. White \emph{can wait} if $z\in L(z)$, and Black if $z\in R(z)$. The move $z\to z$ is a
\emph{pass}: it changes only the side to move. The other options of $Z$ are suicides, the moves of $Z$ that lose. Without
suicides, $Z$ is the D-game form of $\mathrm{on}$ if only White can wait, of $\mathrm{off}$ if only Black can, and of
$\mathrm{dud}$ if both can. For a short D-game $G$ we study $G\boxplus Z$; we write its position $(x,z)$ as $x$. It is not
short when a side can wait.

\textbf{Forcing.} Let $G$ be short. To test forcing from $(x,s)$, let the opponent $\bar\sigma$ choose at each of its
turns any move of $G$ or a pass (whether or not $\bar\sigma$ can wait in $Z$), and let $\sigma$ choose only moves of $G$.
Then $\sigma$ \emph{forces a decisive event against a passing opponent} from $(x,s)$ if $\sigma$ has a strategy under
which every such play reaches the terminal of $\sigma$ after finitely many moves. Since the opponent always has a pass,
$\sigma$ cannot win such a play by leaving it without a move.

\textbf{Predicates.} Let $G$ be short. Going through the positions of $G$ in an order that lists every option that is a
position before the position it is an option of, define for each position $x$,
first $\lambda_{\mathrm L}(x)$ and then $\lambda_{\mathrm R}(x)$:
\begin{itemize}
\item[] $\lambda_{\mathrm L}(x)$: some White option of $x$ is $\oplus$, or is a position $y$ with $\lambda_{\mathrm R}(y)$;
\item[] $\lambda_{\mathrm R}(x)$: $\lambda_{\mathrm L}(x)$ holds, and every Black option of $x$ is $\oplus$ or a
position $y$ with $\lambda_{\mathrm L}(y)$.
\end{itemize}
Define $\varrho_{\mathrm R}(x)$ and $\varrho_{\mathrm L}(x)$ in the same way with the colours exchanged: some Black option
of $x$ is $\ominus$ or a position $y$ with $\varrho_{\mathrm L}(y)$; and $\varrho_{\mathrm R}(x)$ holds and every White
option of $x$ is $\ominus$ or a position $y$ with $\varrho_{\mathrm R}(y)$. The subscript names the side to move
($\mathrm L$ for White, $\mathrm R$ for Black). By Lemma~A.18, $\lambda_s(x)$ says that White forces a decisive event
against a passing opponent from $(x,s)$, and $\varrho_s(x)$ the same for Black; this is the condition for winning when the
opponent can wait (Theorem~A.18a), but not in general. By definition $\lambda_{\mathrm R}(x)$ implies
$\lambda_{\mathrm L}(x)$, and $\varrho_{\mathrm L}(x)$ implies $\varrho_{\mathrm R}(x)$.

\textbf{A.18 Lemma (forcing and blocking).} \emph{Let $G$ be a short D-game, $Z$ as above, $x$ a position of $G$ and $s$
a side to move. (a) If $\lambda_s(x)$, White has a strategy in $G\boxplus Z$ that uses only moves of $G$ and wins
against every play of Black, passes included, whether or not Black can wait; in particular White forces a decisive
event against a passing opponent from $(x,s)$. (b) Let $\lambda_s(x)$ fail, and grant Black a pass at every Black turn
(a pass changes only the side to move). Then Black has a strategy that uses only moves of $G$ and passes, under which no
play reaches $\oplus$, whatever White does (moves of $G$, passes, suicides); in particular White does not force a
decisive event against a passing opponent from $(x,s)$. If Black can wait in $Z$, this is a strategy in $G\boxplus Z$
itself. (c) Parts (a) and (b) also hold with the colours exchanged, $\varrho$ in place of $\lambda$; in (b) White is
granted the pass, and the strategy is one in $G\boxplus Z$ if White can wait.}

\begin{proof}
(a) White keeps the invariant that at the current position $y$, with $t$ to move, $\lambda_t(y)$ holds; this is true at
the start. With White to move, $\lambda_{\mathrm L}(y)$ gives a White option that is $\oplus$, and White wins, or a
position $y'$ with $\lambda_{\mathrm R}(y')$; White plays it. With Black to move, $\lambda_{\mathrm R}(y)$ holds. A Black
pass leads to $y$ with White to move, and $\lambda_{\mathrm L}(y)$ holds since $\lambda_{\mathrm R}$ implies
$\lambda_{\mathrm L}$. A Black move of $G$ leads to $\oplus$ or to a position $y'$ with $\lambda_{\mathrm L}(y')$, by the
second clause of $\lambda_{\mathrm R}(y)$; a Black suicide in $Z$ is $\oplus$. So the invariant holds, no play reaches
$\ominus$, and White always has a move. Every White move is a move of $G$, to $\oplus$ or to a position, and the moves of
$G$ between positions in a play are at most $h-1$, since $G$ is short. White moves at every second turn, so every play
has at most $2h+1$ moves. It ends at $\oplus$ or with Black unable to move; either way White wins. Against a passing
opponent Black always has a move, so the play ends at $\oplus$.

(b) Black keeps the invariant that $\lambda_t(y)$ fails at the current position $y$ with $t$ to move. With Black to move,
$\lambda_{\mathrm R}(y)$ fails. If $\lambda_{\mathrm L}(y)$ fails, Black uses its granted pass. Otherwise the second
clause of $\lambda_{\mathrm R}(y)$ fails, so some Black option of $y$ is $\ominus$, and Black plays it and wins, or is a position
$y'$ where $\lambda_{\mathrm L}(y')$ fails, and Black plays it. With White to move, $\lambda_{\mathrm L}(y)$ fails, so no
White option of $y$ is $\oplus$ and every White option that is a position $y'$ has $\lambda_{\mathrm R}(y')$ false. A
White pass leads to $y$ with Black to move, where $\lambda_{\mathrm R}(y)$ fails because $\lambda_{\mathrm L}(y)$ does; a
White suicide is $\ominus$. So the invariant holds, and neither White nor Black (under this strategy) ever moves to
$\oplus$. Black's prescribed replies are legal in the forcing test with White forcing; every White move of $G$ keeps the
invariant, and if White to move has no option in $G$ the test ends without reaching $\oplus$. So White does not force a
decisive event against a passing opponent. If Black can wait in $Z$, the granted pass is the pass of $Z$, and the strategy is one in
$G\boxplus Z$.

(c) is (a) and (b) with the colours exchanged.
\end{proof}

\textbf{A.18a Theorem (outcomes with spare moves; Theorem~\ref{thm:loopy}).} \emph{Let $G$ be a short D-game, for
instance a sum of short D-games, let $Z$ be as above, and consider $G\boxplus Z$ from a position $x$ with $s$ to move.
\begin{itemize}
\item[(a)] If both sides can wait: White wins if and only if $\lambda_s(x)$, if and only if White forces a decisive event
against a passing opponent; Black wins if and only if $\varrho_s(x)$, if and only if Black forces a decisive event
against a passing opponent; otherwise $(x,s)$ is a draw. A side that wins has a winning strategy that never passes, and
having the move never hurts: $\lambda_{\mathrm R}(x)$ implies $\lambda_{\mathrm L}(x)$, and $\varrho_{\mathrm L}(x)$
implies $\varrho_{\mathrm R}(x)$.
\item[(b)] If only White can wait: every play is finite and $(x,s)$ is not a draw. Black wins if and only if
$\varrho_s(x)$, if and only if Black forces a decisive event against a passing White; otherwise White wins.
\item[(c)] If only Black can wait: the mirror image of (b).
\end{itemize}}

\begin{proof}
(a) If $\lambda_s(x)$, White wins by Lemma~A.18(a), with a strategy that uses only moves of $G$. If $\lambda_s(x)$
fails, Black follows the strategy of Lemma~A.18(b), which is a strategy in $G\boxplus Z$ since Black can wait; under it no
play reaches $\oplus$, and Black always has a move (a pass or the chosen option). A play ends in a White win only at
$\oplus$ or with Black unable to move, so White does not win. The two parts of Lemma~A.18 also show that $\lambda_s(x)$
holds exactly when White forces a decisive event against a passing opponent. The statements for Black follow in the same
way from Lemma~A.18(c), whose blocking strategy for White is a strategy in $G\boxplus Z$ since White can wait. If neither
$\lambda_s(x)$ nor $\varrho_s(x)$ holds, neither side wins, which is a draw: each side, by its blocking strategy, avoids
defeat against every play of the other, and when both follow these strategies the play never ends. The implications are
the remark after the definition of the predicates.

(b) Black cannot pass, so every Black move is a move of $G$ between positions, or a move to a terminal (in $G$ or a
suicide), which ends play. The moves of $G$ between positions in a play, by either side, are at most $h-1$, since passes
leave the position of $G$ unchanged and $G$ is short. So Black makes at most $h$ moves, and since the turns alternate
every play has at most $2h+1$ moves. Each position has finitely many options, so the game from $(x,s)$ is a finite game
tree whose every leaf has a winner (a terminal, or a player with no option), and by Zermelo's theorem one side
wins: there are no draws. If $\varrho_s(x)$, Black wins by Lemma~A.18(c), part (a) with colours exchanged; its strategy uses only moves of $G$ and
beats White's passes. If $\varrho_s(x)$ fails, White follows the strategy of Lemma~A.18(c), part (b) with colours
exchanged, which uses moves of $G$ and passes and is a strategy in $G\boxplus Z$ since White can wait. Under it no play
reaches $\ominus$ and White always has a move; every play is finite, so it ends at $\oplus$ or with Black unable to move, and
White wins. As in (a), $\varrho_s(x)$ holds exactly when Black forces a decisive event against a passing White.

(c) is (b) with the colours exchanged.
\end{proof}

\textbf{A.18b Remark (the statement of Section~\ref{S-kings}).} Theorem~\ref{thm:loopy} is (a) and (b): ``the
components form a short game in which every decisive event ends the game'' is $G=A_1\boxplus\dots\boxplus A_k$ with short
$A_i$; ``let a side that can wait pass at any time'' and ``let every other move in $Z$ lose'' are the D-game $Z$: a side
that can wait may pass at any of its turns, and every other move of $Z$ is a suicide, an immediate win for the opponent.
The forcing and blocking strategies above never prescribe a suicide, and each of them covers an opponent who plays one.
``Against a passing opponent'' means against an opponent who chooses at each turn between a move of $G$ and a pass, not
one who must pass: if $G$ has one position $x$ with $L(x)=\{\oplus\}$ and $R(x)=\{\ominus\}$, and White is to move, Black
reaches $\ominus$ against a White that always passes, yet $\varrho_{\mathrm L}(x)$ fails and White wins by moving to
$\oplus$. In (b) Black can win only at a terminal, since White always has a pass, so ``Black wins'' means ``Black forces $\ominus$''; White may also win because Black runs out of moves. The theorem is about the
D-game. That a king stepping back and forth in chess acts as such a pass is a further hypothesis, not treated here: the
steps must be legal at every position reached, and must change nothing else (in chess a king move also ends an e.p.\
right, while the pass of $Z$ keeps the position of $G$, right included).

\subsubsection*{A.VI A local test for (V2) with frozen kings}

\textbf{Setting.} We use the frozen-king game $\mathrm{F}$ of Part~A.III: chess in which the first promotion ends the game
and wins for the promoter, and a legal king move ends the game at once with a loss for its maker. Checkmate wins; a side
with no legal move (king moves count) that is not in check is stalemated, a draw; infinite play is a draw. The board is
flat or a cylinder with $H\ge5$ ranks. Let $y$ be a legal position (the kings not adjacent, the side not to move not in
check) with pawns and the two kings only, no promotion having occurred, with side $s$ to move and with the
e.p.\ right, if any, created by the double step that led to $y$ (Part~A.I). Fix a side $\beta$, the beneficiary, and let $\omega=\bar\beta$. For a pawn $q$ of either colour let $\rho(q)$ be its rank counted
from $\beta$'s side: $\rho(q)$ is its rank if $\beta$ is White and $H+1$ minus its rank if $\beta$ is Black. So $\beta$'s
pawns move to larger $\rho$, $\omega$'s to smaller $\rho$, and ``ahead of a pawn of $\beta$'' means larger $\rho$. For a
pawn $p$ of $\beta$ let
\[
d=H-\rho(p)\ \text{ if }\rho(p)\ge3,\qquad d=H-3\ \text{ if }\rho(p)=2,
\]
the number of moves $p$ needs to promote when it moves alone, with a double step first if $\rho(p)=2$; and let $e=d$ if
$s=\omega$ and $e=d-1$ if $s=\beta$. If play reaches the $d$-th move of $p$ and $p$ has moved at every turn of $\beta$,
then $\omega$ has made exactly $e$ moves before that move.

\textbf{The game $G_\omega(y;p)$.} Delete $p$ from $y$. Only $\omega$ moves, one move after another. $\beta$'s other
pawns stand still and may be captured; both kings stand still, occupy their squares and are never captured. $\omega$'s
pawns step, double step and capture by the pawn rules of Part~A.I, with the kings as obstacles and with no test of check,
and may capture e.p.\ only as the first move, only if $s=\omega$, and only with the right of $y$ (by (H1) that right
allows no capture of $p$). Let $\mu_\omega(y;p)$ be the least number of moves of a sequence in $G_\omega(y;p)$
after which some pawn of $\omega$ stands on its last rank, or $\infty$ if there is none. This minimum is over all pawns of
$\omega$, passed or not, and over all sequences, captures of $\beta$'s pawns included.

\textbf{Hypotheses.} Let $f$ be the file of $p$; the \emph{strip} is the set of files at distance at most $1$ from $f$.
\begin{itemize}
\item[(H1)] \emph{Clear path.} $p$ passes the e.p.-aware test at $y$ (Part~A.I): no pawn of $\omega$ in the strip has
$\rho>\rho(p)$, and $p$ has not just double stepped with a pawn of $\omega$ beside it. Moreover no pawn and no king of
either colour stands on a square of $f$ ahead of $p$, the promotion square included.
\item[(H2)] \emph{No check.} No pawn of either colour stands on a file at distance $1$ from the file of $\beta$'s king.
\item[(H3)] \emph{No stalemate.} $\omega$'s king has a neighbouring square $t$ at king distance (the maximum of the
file distance and the absolute difference of the ranks) at least $2$ from $\beta$'s king such that no pawn of either colour stands on
the file of $t$ or on a file at distance $1$ from it.
\item[(H4)] \emph{No interference.} Every pawn of $\beta$ other than $p$ on a file at distance $1$ from $f$ has
$\rho\le\rho(p)+1$.
\item[(H5)] \emph{Timing.} $\mu_\omega(y;p)\ge e+1$.
\end{itemize}

\textbf{A.19 Proposition (frozen-king passer test).} \emph{Under (H1)--(H5), $\beta$ wins $\mathrm{F}$ from $y$ with $s$
to move by the strategy $\Pi$: move $p$ forward at every turn, by a double step if $\rho(p)=2$. Every move of $\Pi$ is
legal when it is due, and every play in which $\beta$ follows $\Pi$ ends with a win for $\beta$, at the latest with the
$d$-th move of $p$, which promotes.}

\begin{proof}
Write $\rho_p$ for the current value of $\rho(p)$ and $\rho_0$ for its value at $y$. We consider plays in which $\beta$
follows $\Pi$ and $\omega$ moves arbitrarily, as long as the game has not ended.

\emph{Step 0 (frame).} No promotion has occurred, so the only pieces are pawns and kings, and by Lemma~A.1 a pawn changes
file only by capturing a pawn on the file it enters; a file with no pawn never gets one. A king move ends $\mathrm{F}$, so
in a play that goes on neither king has moved, and both stand where they stood at $y$. Under $\Pi$, $p$ is the only pawn
of $\beta$ that moves, and it stays on $f$.

\emph{Step 1 (checks and stalemate).} By (H2) and Step~0, no pawn ever stands on a file at distance $1$ from the file
of $\beta$'s king, so no pawn of $\omega$ attacks it, and $\omega$'s king, which does not move, is not adjacent to it
since $y$ is legal. So $\beta$ is never in check, and a move of $p$ is illegal only if its squares are occupied. By (H3)
and Step~0, the file of $t$ and the files at distance $1$ from it never hold a pawn; in particular $f$ is not one of them.
So $t$ is never occupied or attacked by a pawn, it is neither the square of $\beta$'s king nor adjacent to it, and the
move of $\omega$'s king to $t$ is legal at every turn of $\omega$, even in check. So $\omega$ is never stalemated.

\emph{Step 2 (invariant).} We show, by induction on the number of moves played, that
\begin{itemize}
\item[] (J$_\beta$) with $\beta$ to move, every pawn of $\omega$ in the strip has $\rho\le\rho_p+1$;
\item[] (J$_\omega$) with $\omega$ to move, every pawn of $\omega$ in the strip has $\rho\le\rho_p$, and $p$ cannot be
captured e.p.;
\item[] (I) no pawn and no king stands on a square of $f$ ahead of $p$.
\end{itemize}
At $y$ these are (H1): with $s=\omega$ it is (J$_\omega$) as stated, and with $s=\beta$ the bound $\rho\le\rho_0$ gives
(J$_\beta$). The kings never move (Step~0), so (I) concerns pawns only after $y$.

\emph{$\beta$ moves.} By (I) the squares of the move of $p$ are empty, and by Step~1 the move is legal; $p$ goes to
$\rho'\in\{\rho_p+1,\rho_p+2\}$ and stays on $f$, so (I) still holds. By (J$_\beta$) every pawn of $\omega$ in the strip
has $\rho\le\rho_p+1\le\rho'$. After a single step there is no right. After a double step, $\rho'=\rho_p+2$, every pawn
of $\omega$ in the strip is below $\rho'$, and a pawn that could capture $p$ e.p.\ would stand beside $p$ at $\rho'$; there
is none. So (J$_\omega$) holds.

\emph{$\omega$ moves a pawn.} A step or a double step lowers $\rho$ and keeps the file, so (J$_\beta$) holds for the
moved pawn. It cannot enter a square of $f$ ahead of $p$: a pawn of $\omega$ on $f$ has $\rho\le\rho_p$ by (J$_\omega$)
and moves to smaller $\rho$. A capture lands on the square of a captured pawn $q$ of $\beta$, or, e.p., on the square that
$q$ skipped. First, $q\ne p$: a capture of $p$ would need a pawn of $\omega$ in the strip at $\rho_p+1$, or a usable
right against $p$, and (J$_\omega$) excludes both. Second, a pawn $q\ne p$ of $\beta$ has not moved since $y$. If $q$ is on
$f$, it stands behind the square of $p$ at $y$ by (H1), and $p$ only advances; so the landing square, the square of $q$
or, e.p., the square behind it, is behind $p$, and (I) is kept. If $q$ is on a file at distance $1$ from $f$, then
$\rho(q)\le\rho_0+1\le\rho_p+1$ by (H4), and the landing rank is $\rho(q)$, or $\rho(q)-1$ for an e.p.\ capture. Captures
landing outside the strip do not concern (J$_\beta$). Third, an e.p.\ capture by $\omega$ needs a right created by
$\beta$; under $\Pi$ only $p$ double steps, and $p$ cannot be captured e.p., so any such capture is made at the first turn
from $y$, with $s=\omega$, using the right of $y$ against a pawn $q\ne p$; it is covered above. So after the move every pawn of
$\omega$ in the strip has $\rho\le\rho_p+1$, which is (J$_\beta$), and (I) holds. (Pawns of $\omega$ never move to
larger $\rho$, but a capture of a pawn of $\beta$ on a file next to $f$ brings a pawn of $\omega$ into the strip at the
rank of the captured pawn, or one rank lower after an e.p.\ capture; (H4) is what bounds it.)

\emph{Step 3 ($p$ survives and is never blocked).} By (J$_\omega$), at no turn of $\omega$ can $p$ be captured, by a
normal or an e.p.\ capture; by (I) its path is empty at every turn of $\beta$. With Step~1, every move of $\Pi$ is legal.

\emph{Step 4 ($\omega$'s pawn moves are moves of $G_\omega$).} Take a prefix of a play that follows $\Pi$ which contains no king
move and no promotion, except that its last move may be a promotion by a pawn of $\omega$. After each move of the prefix, the pieces other than $p$ stand on the same
squares as after the same pawn moves of $\omega$, in the same order, in $G_\omega(y;p)$ (a promoting pawn counted on its arrival square): the kings and $\beta$'s other
pawns have not moved, the same pawns of $\beta$ have been captured in both, and $p$ was never captured (Step~3). We
compare the moves in order. A step or double step of $\omega$ in the play uses empty squares, which are also empty in
$G_\omega$, where $p$ is deleted. A capture takes a pawn $q\ne p$ of $\beta$, which stands on the same square in
$G_\omega$. Any e.p.\ capture by $\omega$ uses the right of $y$ at $\omega$'s first move (Step~2), and the same capture,
by the same pawn, is allowed in $G_\omega$. So the pawn moves of $\omega$, in order, form a sequence of
$G_\omega(y;p)$, and if $\omega$ promotes with its $k$-th move then $k\ge\mu_\omega(y;p)$.

\emph{Step 5 (count).} Let a play follow $\Pi$. Before the $d$-th move of $p$, $\omega$ has at most $e$ turns. If
$\omega$ makes a king move at one of them, $\beta$ wins at once. Otherwise its moves there are pawn moves, and by Step~4
and (H5) none of them promotes, since $e<\mu_\omega(y;p)$. By Step~1, $\omega$ always has a legal move, the king move to
$t$, so it is neither stalemated nor checkmated, and $\beta$ is never in check; by Step~3, $\beta$ always has the move of
$\Pi$. So a play that has not ended with a king move of $\omega$ reaches the $d$-th move of $p$, which is legal by Step~3
and promotes, and $\beta$ wins.
\end{proof}

\textbf{A.19a Corollary ((V2) with frozen kings).} \emph{Let $Q$ be a king group with its context as in Part~A.III, let
$Y\in\mathrm{R}(Q)$ have side $M$ to move, and let $\mu$ be a context move at $Y$ that produces a passed pawn and is
adjudicated for $\sigma$ in $W$. Let $y'$ be the position after $\mu$, with $N=\bar M$ to move. If for some pawn $p$ of
$\sigma$ the hypotheses (H1)--(H5) hold at $y'$ with $s=N$ and $\beta=\sigma$ (so $e=d$ if $\sigma=M$ and $e=d-1$ if
$\sigma=N$), then $\sigma$ wins $y'$ in $\mathrm{F}$. If this holds for every such move at every $Y\in\mathrm{R}(Q)$, then
(V2) holds with chess replaced by $\mathrm{F}$, the form used in Remark~A.12(b).}

\begin{proof}
$y'$ is a legal position reached by a legal move. No promotion has occurred: play from $Q$ to $Y$ contains no decisive
event, hence no Z-move and no context promotion (proof of Theorem~A.11), and $\mu$ does not promote, by Lemma~A.2(c) as
in that proof. So Proposition~A.19 applies at $y'$ with $N$ to move and $\beta=\sigma$.
\end{proof}

\textbf{A.19b Remark (what the test gives and what it does not).} (a) The test is a sufficient condition. For any
fixed race rule, the corollary applies the test to the side that $W$ credits at that move; a different rule can change
the side to be tested. (H5) counts every pawn of $\omega$, passed or not, and every sequence of $\omega$'s moves with $p$
removed, including captures of $\beta$'s other pawns and any e.p.\ capture available at the first move; a count of
passed pawns alone does not suffice. (b) It is a statement about $\mathrm{F}$, where king moves lose; Step~0 uses that the kings never move. It gives
neither $\mathrm{P}$ nor chess. On $\cyl{8}$ with six ranks, with White king a1 and pawn c2, Black king a6 and White to
move, (H1)--(H5) hold for $\beta$ White and $p=$\,c2 ($\mu_\omega=\infty$), so White wins $\mathrm{F}$; yet in
$\mathrm{P}$, with real kings, the position is a draw \lab{C1}. Thus, when the test passes at
every move that Corollary~A.19a quantifies over, the corollary supplies (V2) for Theorem~A.11 in the form of Remark~A.12(b) with $\mathrm{F}$, and not
for $\mathrm{P}$ or chess.
